\section{Estratégia}

\subsection{Oráculo e precedência}

A verificação parte da seguinte precedência:

\begin{center}
documentação normativa $\rightarrow$ CUE $\rightarrow$ Alloy $\rightarrow$
matriz de implementação $\rightarrow$ implementação
\end{center}

Os worklogs M0--M13 registram rationale e contraexemplos, mas não substituem a
redação normativa final. A implementação existente também não altera o contrato
silenciosamente. Um conflito real entre fontes normativas interrompe o trabalho
até que seja relatado e resolvido na fonte competente.

\subsection{Camadas}

\begin{itemize}
  \item testes unitários isolam funções puras ou dependências controladas por mocks;
  \item testes de integração exercitam Admin SDK, callables e serviços locais;
  \item testes de Security Rules usam os Emulators e contextos autenticados;
  \item testes E2E exercitam fluxos observáveis no navegador;
  \item build, typecheck e lint fornecem evidência estrutural, não conformidade funcional.
\end{itemize}

\subsection{TDD incremental e RED/GREEN}

A unidade de trabalho é uma ou poucas features \texttt{IMP-*} fortemente
relacionadas. Primeiro se extrai o contrato e se cria um teste que falha pela
ausência ou divergência pretendida (RED). A menor mudança de implementação é
então feita até o teste passar (GREEN), seguida da regressão relevante. Uma
falha de ambiente, dependência ou comando é \emph{INFRA\_FAILURE}; ela não é um
RED válido de produto.

Uma feature somente pode mudar para \texttt{NÃO DIVERGENTE} depois da revisão
das obrigações normativas e da aprovação das camadas aplicáveis de backend,
Rules, frontend e E2E, sem divergência conhecida restante. Esse estado não
equivale a homologação de produção.

\subsection{Limitações}

Teste aprovado demonstra apenas os cenários cobertos. Mocks não provam integração;
compilação não prova invariantes; e uma suite legada pode codificar contrato
anterior. Resultados devem registrar versão, comando, escopo e classificação da
falha. Ausência de teste não é evidência de conformidade.
